\BourbakiExerciseGroup

\begin{enumerate}
\def\labelenumi{\arabic{enumi})}
\tightlist
\item
  Let \(U = (\alpha_{ij})\) be a matrix with \(m\) rows and \(n\) columns over the principal ideal domain A. a) Suppose first of all that the \(\alpha_{ij}\) are setwise coprime. Show that there exist two invertible square matrices \(P\) and \(Q\) with entries in A, such that one of the entries of the matrix \(PUQ\) is equal to 1. (For each \(\alpha_{ij} \neq 0\) , let \(S(\alpha_{ij})\) be the sum of the exponents in a decomposition of \(\alpha_{ij}\) into irreducible elements, and let \(s(U)\) be the least of the numbers \(s(\alpha_{ij})\) ; if \(s(U) > 0\) , show that there exist two invertible square matrices \(R\) and \(S\) such that
\end{enumerate}

\[
s (R U S) <   s (U);
\]

use Bezout's identity, and restrict the matrices \(R\) and \(S\) considered, after permuting the rows and columns if necessary, to those of the form \(\begin{pmatrix} T & 0 \\ 0 & I \end{pmatrix}\) , where \(T\) is a \(2 \times 2\) matrix and \(I\) an identity matrix, or to products of such matrices and permutation matrices.)

\begin{enumerate}
\def\labelenumi{\alph{enumi})}
\setcounter{enumi}{1}
\tightlist
\item
  If \(\delta_1\) is a gcd of the entries of \(U\) , show that there exist two invertible matrices \(P_1, Q_1\) , such that
\end{enumerate}

\[
P _ {1} U Q _ {1} = \left( \begin{array}{c c} \delta_ {1} & 0 \\ 0 & U _ {1} \end{array} \right)
\]

where all the entries of \(U_{1}\) are divisible by \(\delta_{1}\) (use a)).

\begin{enumerate}
\def\labelenumi{\alph{enumi})}
\setcounter{enumi}{2}
\tightlist
\item
  For \(\mathbf{A} = \mathbf{Z}\) , use \(b\) to obtain a method of calculating the invariant factors of an explicit matrix with entries in \(\mathbf{Z}\) . Apply this method to the matrix
\end{enumerate}

\[
\left( \begin{array}{c c c c} 6 & 8 & 4 & 20 \\ 12 & 12 & 18 & 30 \\ 18 & 4 & 4 & 10 \end{array} \right).
\]

\begin{enumerate}
\def\labelenumi{\arabic{enumi})}
\setcounter{enumi}{1}
\item
  Let A be a principal ideal domain. Then two submodules M and N of \(A^{q}\) are transforms of one another under automorphisms of \(A^{q}\) if and only if they have the same invariant factors with respect to \(A^{q}\) .
\item
  Let A be a principal ideal domain with field of fractions K, let E be a vector space over K, let M be a finitely generated A-module of rank n contained in E, and let N be a finitely generated A-module of rank \(p \leqslant n\) contained in KM. Show that there exists a basis \((e_{i})_{1 \leqslant i \leqslant n}\) of M and a basis of N consisting of vectors \(\alpha_{i}e_{i}\) ( \(1 \leqslant i \leqslant p\) ), where \(\alpha_{i} \in K\) and \(\alpha_{i}\) divides \(\alpha_{i+1}\) (with respect to the ring A) for \(1 \leqslant i \leqslant p - 1\) ; show also that the fractional ideals \(A\alpha_{i}\) are uniquely determined (invariant factors of N with respect to M).
\item
  Let \(G\) be a finite abelian group.
\end{enumerate}

\begin{enumerate}
\def\labelenumi{\alph{enumi})}
\item
  For every prime number \(p\) , show that the number of elements of order \(p\) in \(G\) is \(p^N - 1\) , where \(N = \sum_{n=1}^{\infty} m(p^n)\) (in the notation of Prop. 9 of VII, p.~24).
\item
  Deduce that if the equation \(x^p = 1\) has at most \(p\) solutions in \(G\) , for each prime number \(p\) , then \(G\) is cyclic.
\end{enumerate}

\begin{enumerate}
\def\labelenumi{\arabic{enumi})}
\setcounter{enumi}{4}
\item
  Let \(G\) be a finite abelian group of order \(n\) . Show that for every integer \(q\) dividing \(n\) there exists a subgroup of order \(q\) in \(G\) (decompose \(G\) as a direct sum of indecomposable cyclic subgroups).
\item
  Let A be a commutative ring and let E be an A-module. Let \(\mathcal{E}\) denote the ring \(\operatorname{End}_{\mathbf{z}}(\mathbf{E})\) of endomorphisms of the abelian group E (without operators); for every subring \(\mathbf{B} \subset \mathcal{E}\) , let \(\mathbf{B}'\) denote the centraliser of B in \(\mathcal{E}\) (the subring of \(\mathcal{E}\) consisting of elements which commute with every element of B).
\end{enumerate}

\begin{enumerate}
\def\labelenumi{\alph{enumi})}
\item
  Let \(A_{0}\) be the subring of E consisting of the homotheties \(x \mapsto \lambda x (\lambda \in A)\) ; if a is the annihilator of E, then \(A_{0}\) is known to be isomorphic to A/a, and its centraliser \(A_{0}^{\prime}\) in E is the ring \(\operatorname{End}_{\mathbf{A}}(\mathbf{E})\) of endomorphisms of E as an A-module. Show that the centraliser \(A_{0}^{\prime\prime}\) of \(A_{0}^{\prime}\) is commutative, and that its centraliser \(A_{0}^{\prime\prime\prime}\) is equal to \(A_{0}^{\prime}\) (notice that \(A_{0} \subset A_{0}^{\prime}\) ).
\item
  If the submodule F is a direct factor of E, show that \(u(F) \subset F\) for all \(u \in A_{0}''\) (consider the projection of E onto F obtained from the decomposition of E as a direct sum of F and another submodule).
\item
  Suppose that E is the direct sum of a family \((F_{i})\) of cyclic submodules. If \(u \in A_{0}''\) show that, for each index i, there exists an element \(\alpha_{i} \in A\) such that \(u(x) = \alpha_{i}x\) for all \(x \in F_{i}\) (use b)).
\item
  Suppose E is the direct sum of a sequence \((\mathrm{E}_{n})\) (finite or infinite) of cyclic submodules such that, if \(b_{n}\) denotes the annihilator of \(E_{n}\) , then the sequence \((\mathfrak{b}_{n})\) is increasing. Show that \(A_{0}^{\prime\prime}=A_{0}\) (use c), and the fact that for each index i\textgreater1 there exists an endomorphism \(v_{i}\in A_{0}^{\prime}\) which maps \(E_{i-1}\) into \(E_{i}\) .
\end{enumerate}

Apply this to the case where A is a principal ideal domain and E is a finitely generated A-module.

\begin{enumerate}
\def\labelenumi{\arabic{enumi})}
\setcounter{enumi}{6}
\tightlist
\item
  Let A be a principal ideal domain with field of fractions K, and let E be a finitely generated, torsion-free A-module of rank n.~Let \(E^{*}\) be the dual of E; then E and \(E^{*}\) are each isomorphic to \(A^{n}\) .
\end{enumerate}

\begin{enumerate}
\def\labelenumi{\alph{enumi})}
\item
  Let M be a submodule of E; show that, if \(M^{0}\) is the submodule of \(E^{*}\) orthogonal to M, then the module \(E^{*}/M^{0}\) is torsion-free, and consequently \(M^{0}\) is a direct factor in \(E^{*}\) ; the submodule \(M^{00}\) of E orthogonal to \(M^{0}\) is equal to \(E \cap KM\) (considering E as being canonically embedded in an n-dimensional vector space over K). We can consider \(E^{*}/M^{0}\) as being canonically embedded in the dual \(M^{*}\) of M; show that the invariant factors of \(E^{*}/M^{0}\) with respect to \(M^{*}\) are equal to the invariant factors of M with respect to E (use Th. 1 of VII, p.~18).
\item
  Let F be a second finitely generated, torsion-free A-module, and let u be a linear transformation from E into F. Show that the invariant factors of \(u(F^{*})\) with respect to \(E^{*}\) are equal to those of \(u(E)\) with respect to F (proceed as in Prop. 5 of VII, p.~21).
\end{enumerate}

\begin{enumerate}
\def\labelenumi{\arabic{enumi})}
\setcounter{enumi}{7}
\tightlist
\item
  Let G be a module of finite length over a principal ideal domain A. For any pair of free modules M and N such that \(N \subset M\) and G is isomorphic to M/N, the invariant factors of N with respect to M which are distinct from A are identical to the invariant factors of G (VII, p.~20, Def. 2); their number is called the rank of G.
\end{enumerate}

\begin{enumerate}
\def\labelenumi{\alph{enumi})}
\item
  Show that the rank r of G is the least number of cyclic submodules of G of which G is the direct sum, and is equal to the greatest of the ranks of the \(\pi\) -primary components of G. b) Show that the rank of a submodule or a quotient module of G is at most equal to the rank of G (consider G as a quotient of two free modules of rank r).
\item
  For each \(\lambda \in \mathbf{A}\) , show that the rank of the submodule \(\lambda G\) is equal to the number of invariant factors of \(G\) which do not divide \(\lambda\) (notice that, if \(E = A / (A\alpha)\) , then the module \(\lambda E\) is isomorphic to \((A\lambda) / ((A\alpha) \cap (A\lambda))\) . Deduce that the \(k\) -th invariant factor of \(G\) (in decreasing order) is a gcd of those \(\lambda \in A\) such that \(\lambda G\) has rank \(\leqslant r - k\) .
\item
  Let \(A\alpha_{k}\) ( \(1 \leqslant k \leqslant r\) ) be the invariant factors of G arranged in decreasing order. Let H be a submodule of G of rank r - q, and let \(A\beta_{k}\) ( \(1 \leqslant k \leqslant r - q\) ) be its invariant factors, arranged in decreasing order. Deduce from b) and c) that \(\beta_{k}\) divides \(\alpha_{k+q}\) for \(1 \leqslant k \leqslant r - q\) . Similarly show that, if G/H has rank r - p and invariant factors \(A\gamma_{k}\) ( \(1 \leqslant k \leqslant r - p\) ) in decreasing order, then \(\gamma_{k}\) divides \(\alpha_{k+p}\) for \(1 \leqslant k \leqslant r - p\) .
\item
  Conversely, let L be a torsion module of rank r - q, whose invariant factors \(A\lambda_{k}\) ( \(1 \leqslant k \leqslant r - q\) ) are such that \(\lambda_{k}\) divides \(\alpha_{k+q}\) for \(1 \leqslant k \leqslant r - q\) . Show that there exist submodules M and N of G such that L is isomorphic to M and to G/N (decompose G into a direct sum of submodules isomorphic to A/A \(\alpha_{k}\) ).
\end{enumerate}

\begin{enumerate}
\def\labelenumi{\arabic{enumi})}
\setcounter{enumi}{8}
\tightlist
\item
  Let A be a principal ideal domain and K its field of fractions; let E be a vector space over K, let M be a finitely generated A-module of rank n contained in E, let N be a finitely generated A-module of rank p contained in KM, and let P be a finitely generated A-module of rank q contained in KN.
\end{enumerate}

\begin{enumerate}
\def\labelenumi{\alph{enumi})}
\item
  Let \(\mathbf{A}\alpha_{i}\) ( \(1 \leqslant i \leqslant p\) ) be the invariant factors of \(\mathbf{N}\) with respect to \(\mathbf{M}\) , arranged in decreasing order (Ex. 3). Show that, for \(1 \leqslant k \leqslant p\) , the element \(\alpha_{k}\) is a gcd of those elements \(\lambda \in \mathbf{K}\) such that the rank of the quotient module \((\mathbf{N} + \lambda (\mathbf{KN} \cap \mathbf{M})) / \mathbf{N}\) is \(\leqslant p - k\) (cf.~Ex. 8, d)).
\item
  Let \(\mathbf{A}\beta_{j}\) ( \(1 \leqslant j \leqslant q\) ) be the invariant factors of \(\mathbf{P}\) with respect to \(\mathbf{M}\) and \(\mathbf{A}\gamma_{j}\) ( \(1 \leqslant j \leqslant q\) ) the invariant factors of \(\mathbf{P}\) with respect to \(\mathbf{N}\) , arranged in decreasing order. Show that \(\alpha_{1}\gamma_{j}\) divides \(\beta_{j}\) for \(1 \leqslant j \leqslant q\) (cf.~Ex. 8).
\item
  Show that \(\gamma_{1}\alpha_{j}\) divides \(\beta_{j}\) for \(1 \leqslant j \leqslant q\) . (First consider the case \(n = p = q\) , and apply Ex. 8. In general, show that we may always assume \(p = n\) , and consider a submodule Q of M such that \(\mathbf{P} + \rho \mathbf{Q}\) has rank \(p\) for all \(\rho \in \mathbf{A}\) ; then choose the ideal \(\mathbf{A}\rho\) sufficiently small, and apply Prop. 4 of VII, p.~20.)
\end{enumerate}

\(d\) ) For any submodule H of N of rank \(k \leqslant p\) , let \((\mu_{\mathrm{H}})\) be the fractional ideal consisting of \(\mu \in \mathbf{K}\) such that \(\mu (\mathbf{M} \cap \mathbf{KH}) \subset \mathbf{H}\) . Deduce from \(c\) ) that for each \(1 \leqslant k \leqslant p\) the element \(\alpha_{k}\) is a gcd of the \(\mu_{\mathrm{H}}\) as H runs through all submodules of N of rank \(k\) .

\begin{enumerate}
\def\labelenumi{\arabic{enumi})}
\setcounter{enumi}{9}
\tightlist
\item
  Let A be a principal ideal domain and let M be a submodule of rank n in \(E = A^{n}\) ; let \(A\alpha_{i} (1 \leqslant i \leqslant n)\) be the invariant factors of M with respect to E, arranged in decreasing order. Let N be a submodule of M admitting a complement in M (that is, such that \(N = M \cap KN\) , where K is the field of fractions of A).
\end{enumerate}

\begin{enumerate}
\def\labelenumi{\alph{enumi})}
\item
  Let P be a complement of N in M. Show that the submodule \((P + (E \cap KN))/M\) of E/M is isomorphic to \((E \cap KN)/N\) . Deduce that, if N has rank p and \(A\beta_{k}\) ( \(1 \leqslant k \leqslant p\) ) are the invariant factors of N with respect to E, then \(\beta_{k}\) divides \(\alpha_{k+n-p}\) for \(1 \leqslant k \leqslant p\) (use Ex. 8, d)).
\item
  Give an example to show that the module E/M is not necessarily isomorphic to the product of the modules \((E \cap KN)/N\) and \(E/(P + (E \cap KN))\) (take E/M to be indecomposable).
\end{enumerate}

\begin{enumerate}
\def\labelenumi{\arabic{enumi})}
\setcounter{enumi}{10}
\item
  Let A be a principal ideal domain, let H be a submodule of \(E = A^{n}\) , and let \(H^{0}\) be the submodule of the dual \(E^{*}\) of E orthogonal to H (Ex. 7). Let M be a submodule of rank p of E and let \(A\alpha_{i}\) ( \(1 \leqslant i \leqslant p\) ) be the invariant factors of M with respect to E, arranged in decreasing order. Let \((\nu_{\mathrm{H}})\) be the gcd of the ideals \(f(\mathbf{M})\) of A as f runs through \(H^{0}\) ; show that, for \(1 \leqslant k \leqslant p\) , the element \(\alpha_{k}\) is a gcd of the \(\nu_{H}\) for all submodules H of E of rank k - 1 which admit a complement in E (first consider the case k = 1; apply the result to the submodule \((\mathbf{M} + \mathbf{H})/\mathbf{H}\) of E/H, and use Ex. 8, d), applied to the quotient module \(E/(M + H)\) of E/M).
\item
  Let G be a module of finite length over a ring A. Show that a submodule H of G can only be isomorphic to G if it is equal to G (cf.~II, p.~212, Prop. 16).
\item
  Let A be a principal ideal domain with field of fractions K.
\end{enumerate}

\begin{enumerate}
\def\labelenumi{\alph{enumi})}
\item
  Show that, for every A-module M, the canonical homomorphism \(c_{\mathbf{M}}: \mathbf{M} \to \mathbf{D}(\mathbf{D}(\mathbf{M}))\) is injective (cf.~II, p.~389, Ex. 13).
\item
  If \(u: \mathbf{M} \to \mathbf{N}\) is a homomorphism of A-modules, define natural isomorphisms from \(D(\operatorname{Im}(u))\) onto \(\operatorname{Im}(D(u))\) , from \(D(\operatorname{Ker}(u))\) onto \(\operatorname{Coker}(D(u))\) , and from \(D(\operatorname{Coker}(u))\) onto \(\operatorname{Ker}(D(u))\) .
\item
  Let M be an A-module and N a submodule of M. In the notation of Sect. 9, show that \(N^{00} = N\) (use a) and b)). The module \(N^{0}\) has finite length if and only if M/N has finite length. Conversely, for any submodule \(N'\) of D(M) of finite length, the module \(N'^{0}\) is a submodule of M such that \(M/N^{0}\) has finite length, and \(N' = N'^{00}\) .
\end{enumerate}

\(d\) ) If \(\mathbf{M}\) is a torsion A-module, and \(\mathbf{M}_{\pi}\) denotes the \(\pi\) -primary component of \(\mathbf{M}\) , show that \(\mathrm{D}(\mathbf{M})\) is isomorphic to \(\prod_{\pi \in \mathbb{P}}\mathrm{Hom}(\mathbf{M}_{\pi},\mathbf{U}_{\pi})\) , where \(\mathbf{P}\) is a system of representatives of irreducible elements of A and \(U_{\pi}\) is the \(\pi\) -primary component of K/A for each \(\pi \in P\) (VII, p.~54 Ex. 3).

\begin{enumerate}
\def\labelenumi{\alph{enumi})}
\setcounter{enumi}{4}
\tightlist
\item
  Deduce from d) that \(\mathrm{D}(\mathrm{K}/\mathrm{A})\) is isomorphic to the product \(\prod_{\pi\in P}\hat{\mathbf{A}}_{(\pi)}\) , where
\end{enumerate}

\(\hat{\mathbf{A}}_{(\pi)}\) denotes the completion of the ring \(\mathbf{A}_{(\pi)}\) defined in VII, p.~54, Ex. 2, with respect to the topology defined by taking the ideals \(\pi^n\mathbf{A}_{(\pi)}\) to be a basis of open neighbourhoods of 0 (this completion can also be identified with the inverse limit \(\lim (\mathbf{A} / \mathbf{A}\pi^n)\) ).

\begin{enumerate}
\def\labelenumi{\alph{enumi})}
\setcounter{enumi}{5}
\item
  If M is an A-module of finite length, and a is its annihilator, show that D(M) is isomorphic to the dual of the faithful (A/a)-module associated to M.
\item
  Let M be an A-module of finite length. Show that if N is a submodule of M, then there exists a submodule P of M such that M/N is isomorphic to P and M/P is isomorphic to N (use VII, p.~26, Prop. 10 and VII, p.~27, Prop. 12); two submodules of M with these properties are called reciprocal. Show that, for all \(\alpha \in A\) , the submodule \(M(\alpha)\) of M consisting of all \(x \in M\) such that \(\alpha x = 0\) , and the submodule \(\alpha M\) , are reciprocal.
\end{enumerate}

¶ 14) Let M be a module of finite length over a principal ideal domain A with field of fractions K.

\begin{enumerate}
\def\labelenumi{\alph{enumi})}
\tightlist
\item
  Show that if N is a submodule of M then the rank of M is at most equal to the sum of the ranks of N and M/N (let M = E/H, where E is a free module of rank n and H is a submodule of rank n of E; if N = L/H, where H ⊆ L ⊆ E, notice that, if the rank of N is p, then there exists a submodule R of H, of rank n - p, admitting a complement S in L such that S ∩ H is a complement of R in H; moreover, (E ∩ KR)/R has rank ≥ n - p and is isomorphic to a quotient module of E/L (Ex. 10, a)); complete the proof using Ex. 8, b)). b) Let N be a submodule of M of rank p, and let q be the rank of M/N. Let Aα \(_{i}\) (1 ≤ i ≤ n) be the invariant factors of M and Aβ \(_{k}\) (1 ≤ k ≤ p) those of N, arranged in decreasing order. Show that β \(_{k}\) is a multiple of α \(_{k- (p + q - n)}\) (use a) and Ex. 8, c), noticing that (λM)/(λN) is isomorphic to a quotient module of M/N).
\end{enumerate}

\begin{enumerate}
\def\labelenumi{\arabic{enumi})}
\setcounter{enumi}{14}
\tightlist
\item
  \begin{enumerate}
  \def\labelenumii{\alph{enumii})}
  \tightlist
  \item
    Let M be a module of finite length and rank n over a principal ideal domain A, and let \(\mathbf{A}\alpha_{i}\) ( \(1 \leqslant i \leqslant n\) ) be its invariant factors in decreasing order. Let \((\beta_{i})\) ( \(1 \leqslant i \leqslant n\) ) be a sequence of elements of A such that \(\beta_{i}\) divides \(\alpha_{i}\) for \(1 \leqslant i \leqslant n\) and \(\alpha_{i}\) divides \(\beta_{i+1}\) for \(1 \leqslant i \leqslant n-1\) . Let \((\mathbf{M}_{i})_{1 \leqslant i \leqslant n}\) be a sequence of submodules of M such that M is the direct sum of the \(M_{i}\) and each \(M_{i}\) is isomorphic to \(A/A\alpha_{i}\) ; let \(a_{i}\) be a generator of \(M_{i}\) for \(1 \leqslant i \leqslant n\) . Put
  \end{enumerate}
\end{enumerate}

\[
b = \beta_ {1} a _ {1} + \frac {\beta_ {1} \beta_ {2}}{\alpha_ {1}} a _ {2} + \frac {\beta_ {1} \beta_ {2} \beta_ {3}}{\alpha_ {1} \alpha_ {2}} a _ {3} + \dots + \frac {\beta_ {1} \beta_ {2} \dots \beta_ {n}}{\alpha_ {1} \alpha_ {2} \dots \alpha_ {n - 1}} a _ {n}.
\]

Show that the quotient of M by the cyclic submodule Ab is isomorphic to the direct sum of the n modules \(A/A\beta_{i}\) ( \(1 \leq i \leq n\) ).

\begin{enumerate}
\def\labelenumi{\alph{enumi})}
\setcounter{enumi}{1}
\tightlist
\item
  Let M and N be two modules of finite length over A; let \(A\alpha_{i}\) ( \(1 \leqslant i \leqslant n\) ) and \(A\beta_{k}\) ( \(1 \leqslant k \leqslant p\) ) be the invariant factors of M and N respectively, arranged in decreasing order; suppose that \(p \leqslant n\) , that \(\beta_{k}\) divides \(\alpha_{k+n-p}\) for \(1 \leqslant k \leqslant p\) , and that \(\alpha_{k}\) divides \(\beta_{k+(p+q-n)}\) for \(1 \leqslant k \leqslant n-q\) , where q is an integer such that
\end{enumerate}

\[
q \leqslant n \leqslant p + q.
\]

Show that there exists a submodule P of M, isomorphic to N, such that M/P has rank \(\leq q\) (decompose each of the two modules M and N as a direct sum of q modules, in such a way as to reduce the problem to the case q = 1; then use a) and Ex. 13, g)).

\begin{enumerate}
\def\labelenumi{\alph{enumi})}
\setcounter{enumi}{2}
\tightlist
\item
  For a module N over A to be isomorphic to a direct factor of a module M of finite length, it is necessary and sufficient that every elementary divisor of N be an elementary divisor of M, and that its multiplicity in N be at most equal to its multiplicity in M. Hence obtain an example of a module M of finite length and a submodule N of M such that the rank of M is equal to the sum of the ranks of N and M/N, and such that N does not admit a complement in M (use a)).
\end{enumerate}

\begin{enumerate}
\def\labelenumi{\arabic{enumi})}
\setcounter{enumi}{15}
\tightlist
\item
  Let M be a module over a commutative ring A. A submodule N of M is called characteristic if \(u(N) \subset N\) for every endomorphism u of M. For all \(\alpha \in A\) , the submodule \(M(\alpha)\) of elements annihilated by \(\alpha\) , and the submodule \(\alpha M\) , are characteristic.
\end{enumerate}

\begin{enumerate}
\def\labelenumi{\alph{enumi})}
\item
  If N is a characteristic submodule of M, and P and Q are two complementary submodules of M, show that N is the direct sum of \(N \cap P\) and \(N \cap Q\) (consider the projections from M onto P and Q).
\item
  Let M be a module of finite length over a principal ideal domain A; let \(A\alpha_{i}\) be the invariant factors of M, arranged in decreasing order \((1 \leqslant i \leqslant n)\) ; let \((\mathbf{M}_{i})_{1 \leqslant i \leqslant n}\) be a sequence of submodules of M such that M is the direct sum of the \(M_{i}\) and each \(M_{i}\) is isomorphic to \(A/A\alpha_{i}\) . Let N be a characteristic submodule of M and let \(A\beta_{i}\) be the annihilator of \(N \cap M_{i}\) \((1 \leqslant i \leqslant n)\) ; if \(\alpha_{i} = \beta_{i}\gamma_{i}\) , show that \(\beta_{i}\) divides \(\beta_{i+1}\) and \(\gamma_{i}\) divides \(\gamma_{i+1}\) for all \(1 \leqslant i \leqslant n-1\) (notice that, for \(1 \leqslant i \leqslant n-1\) , there exists an endomorphism of M which maps \(M_{i+1}\) onto \(M_{i}\) , and an endomorphism of M which maps \(M_{i}\) onto a submodule of \(M_{i+1}\) ).
\item
  Conversely, let N be a submodule of M which is a direct sum of submodules \(N_{i} \subset M_{i}\) ; suppose the annihilators \(A\beta_{i}\) of the \(N_{i}\) ( \(1 \leqslant i \leqslant n\) ) satisfy the above conditions. Then show that N is a characteristic submodule of M. Deduce that if two characteristic submodules of M have the same invariant factors then they are equal.
\end{enumerate}

\begin{enumerate}
\def\labelenumi{\arabic{enumi})}
\setcounter{enumi}{16}
\tightlist
\item
  Let A be a principal ideal domain, let \(\alpha\) be a nonzero element of A, and let E be the module A/A \(\alpha\) . Then the product \(E^{n}\) is isomorphic to \(A^{n}/(\alpha A^{n})\) .
\end{enumerate}

\begin{enumerate}
\def\labelenumi{\alph{enumi})}
\item
  Let \(u\) be a linear map from \(\mathbf{E}^n\) into \(\mathbf{E}^m\) ; show that \(u\) can be obtained from a linear map \(v\) from \(\mathbf{A}^n\) into \(\mathbf{A}^m\) on passing to quotients.
\item
  Let \(A\beta_{i}\) ( \(1 \leqslant i \leqslant p \leqslant \inf(m, n)\) ) be the invariant factors of \(v(A^{n})\) with respect to \(A^{m}\) , arranged in decreasing order; let \(q \leqslant p\) be the greatest of the indices i such that \(\alpha\) does not divide \(\beta_{i}\) , and let \(\delta_{i}\) be a gcd of \(\alpha\) and \(\beta_{i}\) for \(1 \leqslant i \leqslant q\) . Show that the kernel \(u^{-1}(0)\) of u is the direct sum of n - q modules isomorphic to E and q cyclic modules isomorphic to \(A/A\delta_{i}\) respectively ( \(1 \leqslant i \leqslant q\) ). If \(\alpha = \gamma_{i}\delta_{i}\) ( \(1 \leqslant i \leqslant q\) ), show that the module \(u(E^{n})\) has invariant factors \(A\gamma_{i}\) .
\end{enumerate}

\begin{enumerate}
\def\labelenumi{\arabic{enumi})}
\setcounter{enumi}{17}
\tightlist
\item
  Let \(\mathbf{C}\) be a principal ideal domain. Consider a system of linear equations
\end{enumerate}

\[
\sum_ {j = 1} ^ {n} \alpha_ {i j} \xi_ {j} = \beta_ {i} (1 \leqslant i \leqslant m)
\]

where the coefficients and the right hand sides belong to C, and let A denote the matrix \((\alpha_{ij})\) with m rows and n columns, and B the matrix obtained by extending A by the \((n+1)\) -st column \((\beta_{i})\) . Show that the system admits at least one solution in \(C^{n}\) if and only if the following conditions are satisfied: 1) A and B have the same rank p; 2) a gcd of the minors of order p of A is equal to a gcd of the minors of order p of B (use Prop. 4 of VII, p.~20 and Ex. 9, c)).

\begin{enumerate}
\def\labelenumi{\arabic{enumi})}
\setcounter{enumi}{18}
\tightlist
\item
  Let \(\mathbf{C}\) be a principal ideal domain. Consider a system of \(m\) « linear congruences » \(\sum_{j=1}^{n} \alpha_{ij} \xi_j \equiv \beta_i (\bmod \alpha)\) , where the coefficients, the right hand sides, and the element \(\alpha\) all belong to C, and where \(\alpha\) is nonzero; let \(A\) denote the matrix \((\alpha_{ij})\) , let \(B\) denote the matrix obtained by extending A by the column \((\beta_i)\) , and let \(\delta_k\) (resp. \(\delta_k'\) ) be a gcd of the \(k\) -th order minors of \(A\) (resp. \(B\) ). Then the system admits at least one solution in \(\mathbf{C}^n\) if and only if: \(a)\) If \(m \leqslant n\) , a gcd of \(\alpha^m\) , \(\alpha^{m-1}\delta_1\) , \ldots, \(\alpha\delta_{m-1}\) , \(\delta_m\) is equal to a gcd of \(\alpha^m\) , \(\alpha^{m-1}\delta_1'\) , \ldots, \(\alpha\delta_{m-1}'\) , \(\delta_m'\) ;
\end{enumerate}

\begin{enumerate}
\def\labelenumi{\alph{enumi})}
\setcounter{enumi}{1}
\tightlist
\item
  if \(m > n\) , a gcd of \(\alpha^{n+1}\) , \(\alpha^n \delta_1, \ldots, \alpha \delta_n\) is equal to a gcd of \(\alpha^{n+1}\) , \(\alpha^n \delta_1', \ldots, \delta_{n+1}'\) .
\end{enumerate}

(Reduce the system of linear congruences to a system of linear equations, and use Ex. 18.)

\begin{enumerate}
\def\labelenumi{\arabic{enumi})}
\setcounter{enumi}{19}
\tightlist
\item
  \begin{enumerate}
  \def\labelenumii{\alph{enumii})}
  \tightlist
  \item
    Let \(\mathbf{C}\) be a principal ideal domain, and let
  \end{enumerate}
\end{enumerate}

\[
f (x _ {1},.., x _ {p}) = \sum_ {(i _ {k})} \alpha_ {i _ {1} \dots i _ {p}} \xi_ {1, i _ {1}} \dots \xi_ {p, i _ {p}}
\]

be a \(p\) -linear form on \(\mathbf{E}^p\) , where \(\mathbf{E}\) is the C-module \(\mathbf{C}^n\) (and where \(\xi_{ij}\) denotes the \(j\) -th coordinate of \(x_i\) ). Show that if \(\delta\) is a gcd of the coefficients \(\alpha_{i_1\dots i_p}\) then there exists an element \((a_k)\in \mathbf{E}^p\) such that \(f(a_1,\dots,a_p) = \delta\) (reduce to the case \(\delta = 1\) , and argue by induction on \(p\) , using Th. 1 of VII, p.~18).

\begin{enumerate}
\def\labelenumi{\alph{enumi})}
\setcounter{enumi}{1}
\tightlist
\item
  Let n \textgreater{} 1 be an integer, and \(\alpha_{i}\) ( \(1 \leqslant i \leqslant n\) ) be n elements of C with gcd \(\delta\) . Show that there exists a square matrix A of order n over C whose first column is \((\alpha_{i})\) and whose determinant is \(\delta\) (use a)).
\end{enumerate}

\begin{enumerate}
\def\labelenumi{\arabic{enumi})}
\setcounter{enumi}{20}
\tightlist
\item
  Let \(C\) be a principal ideal domain, and let \(A = (\alpha_{ij})\) be a square matrix of order \(n\) over \(C\) .
\end{enumerate}

\begin{enumerate}
\def\labelenumi{\alph{enumi})}
\item
  Show that there exists a square matrix U of order n over C, with determinant 1, such that the last column of the matrix UA has its first \((n-1)\) entries zero, and its last entry equal to the gcd of the entries in the n-th column of A (use Ex. 20, b)).
\item
  Deduce from a) that there exists a square matrix V of order n over C, with determinant 1, such that the matrix \(VA = (\beta_{ij})\) is lower triangular (« Hermite reduced form of A »).
\item
  If W is a square matrix of order n over C, with determinant 1, such that \(WA = (\gamma_{ij})\) is lower triangular, show that \(\beta_{ii}\) and \(\gamma_{ii}\) are associate for \(1 \leqslant i \leqslant n\) .
\end{enumerate}

\begin{enumerate}
\def\labelenumi{\arabic{enumi})}
\setcounter{enumi}{21}
\tightlist
\item
  A module M over a principal ideal domain is said to have finite rank if there exists an integer \(m \geqslant 0\) with the following property: every finitely generated submodule of M is contained in a submodule generated by m elements or fewer. The smallest integer m with this property is called the rank of M.
\end{enumerate}

\begin{enumerate}
\def\labelenumi{\alph{enumi})}
\item
  Show that if \(\mathbf{M}\) has finite rank, then every quotient module has finite rank, at most equal to that of \(\mathbf{M}\) .
\item
  For a torsion-free module M, the notion of rank defined above coincides with that defined in II, p.~315.
\item
  Show that if \(\mathbf{M}\) has finite rank then every submodule of \(\mathbf{M}\) has finite rank, at most equal to that of \(\mathbf{M}\) .
\item
  Let \(\mathbf{M}\) be a \(\pi\) -primary module such that \(\mathbf{M}(\pi) = \mathbf{M}\) (in the notation of VII, p.~7). Show that if \(\mathbf{M}\) has finite rank then it is finitely generated.
\item
  Let M be a \(\pi\) -primary module with no element of infinite height. Show that if M has finite rank then it is finitely generated, and the notion of rank coincides with that defined in
\end{enumerate}

Ex. 8. (Notice that if \(\mathbf{M}(\pi)\) is finitely generated then so is \(\mathbf{M}\) , by proving that the heights of elements of \(\mathbf{M}(\pi)\) in \(\mathbf{M}\) are bounded; use Ex. 5 of VII, p.~55 to show this.) \(f)\) Let \(\mathbf{M}\) be a \(\pi\) -primary module of finite rank \(r\) , and let \(\mathbf{M}_0\) denote the submodule consisting of all the elements of infinite height in \(\mathbf{M}\) . Show that \(\mathbf{M}_0\) is a divisible module, and the direct sum of a finite number \(p \leqslant r\) of indecomposable submodules \(\mathrm{U}_{\pi}\) (VII, p.~54, Ex. 3), and that \(\mathbf{M}\) is the direct sum of \(\mathbf{M}_0\) and a finitely generated submodule \(\mathbf{N}\) of rank \(r - p\) . (Notice that \(\mathbf{M} / \mathbf{M}_0\) is finitely generated, using \(e\) ); then apply Ex. 3 of VII, p.~54.) \(g)\) Deduce from the above that an A-module \(\mathbf{M}\) has finite rank \(r\) if and only if the quotient module \(\mathbf{M} / \mathbf{T}\) of \(\mathbf{M}\) by its torsion submodule \(\mathbf{T}\) has finite rank \(p \leqslant r\) , and that each \(\pi\) -primary component of \(\mathbf{T}\) has rank \(\leqslant r - p\) , with at least one of them having rank \(r - p\) .

\begin{enumerate}
\def\labelenumi{\alph{enumi})}
\setcounter{enumi}{7}
\item
  Let \((p_{n})\) be the sequence of distinct prime numbers; let E be the direct sum of the cyclic Z-modules \(Z/p_{n}^{2}Z\) , and let \(e_{n}\) denote the element whose coordinates all vanish except the nth, which is equal to the residue class of \(1 \mod p_{n}^{2}\) . Let M be the submodule of the Z-module \(E \times Q\) generated by the elements \((e_{n}, 1/p_{n})\) . Show that M is a Z-module of rank 2, that the torsion submodule T of M is generated by the elements \((p_{n}e_{n}, 0)\) , and that T is not a direct factor of M. (If \(\bar{x}_{n}\) is the class of \((e_{n}, 1/p_{n}) \mod T\) , show that for all \(m \neq n\) there exists a class \(\bar{y}_{nm} \mod T\) such that \(\bar{x}_{n} = p_{m}\bar{y}_{nm}\) , but that there is no element \(x_{n}\) in the class \(\bar{x}_{n}\) such that for all \(m \neq n\) there exists \(y_{nm} \in M\) with \(x_{n} = p_{m}y_{nm}\) .)
\item
  If \(\mathbf{M}\) is a finitely generated A-module, show that the rank of \(\mathbf{M}\) is equal to the smallest cardinal \(\gamma (\mathbf{M})\) of any generating set for \(\mathbf{M}\) .
\item
  Let M be a torsion A-module of finite rank. Show that if N is a proper submodule of M then N cannot be isomorphic to M (use f).
\end{enumerate}

\begin{enumerate}
\def\labelenumi{\arabic{enumi})}
\setcounter{enumi}{22}
\item
  Let A be a principal ideal domain and let M be a finitely generated A-module. Show that if N and P are two A-modules such that \(M \oplus N\) and \(M \oplus P\) are isomorphic, then N and P are isomorphic. (Reduce to the case where M is cyclic; by identifying \(M \oplus N\) with \(M \oplus P\) , we may consider an A-module E decomposed in two ways as a direct sum of submodules \(F \oplus N\) and \(G \oplus P\) , where F and G are cyclic and isomorphic. In the case where F and G are isomorphic to A, show that if \(N \neq P\) then N is the direct sum of \(N \cap P\) and a submodule isomorphic to A. If F and G are isomorphic to \(A/A\pi^{n}\) , where \(\pi\) is an irreducible element of A, consider two generators a, b of F, G respectively, and examine the cases \(\pi^{n-1}a \notin G\) , \(\pi^{n-1}b \notin F\) , and finally the case where both \(\pi^{n-1}a \in G\) and \(\pi^{n-1}b \in F\) ; in this last case, consider the element \(a + b\) of E and the submodule it generates.)
\item
  An abelian group G is isomorphic to a subgroup of the multiplicative group \(K^{*}\) of nonzero elements of some commutative field K if and only if the torsion subgroup T of G has rank \(\leqslant 1\) (Ex. 22). (For the necessity of this condition, cf.~V, p.~79, Prop. 2; use the same reference to show that the condition is sufficient for T to be isomorphic to a subgroup of \(E^{*}\) , for some field E. For all \(\alpha \in G/T\) let \(g_{\alpha}\) be a coset representative of \(\alpha\) in G, then, identifying T with a subgroup of \(E^{*}\) , we may write \(g_{\alpha}g_{\beta} = e_{\alpha\beta}g_{(\alpha\beta)}\) for all \(\alpha\) , \(\beta \in G/T\) , where \(e_{\alpha\beta} \in E\) . Show that there exists an E-algebra B with basis \((b_{\alpha})_{\alpha \in G/T}\) such that \(b_{\alpha}b_{\beta} = e_{\alpha\beta}b_{(\alpha\beta)}\) for all \(\alpha\) , \(\beta\) , and prove that B is an integral domain using VI, p.~33, Ex. 20, c). Then the field of fractions K of B satisfies the desired conditions.)
\item
  Generalise Th. 1 of VII, p.~18 to left modules over a ring A which is not necessarily commutative, in which every left ideal and every right ideal is principal (an example of such a ring is the tensor product \(\mathbf{K}[\mathbf{X}] \otimes_{\mathbf{K}} \mathbf{D}\) , where \(\mathbf{K}\) is a commutative field and \(\mathbf{D}\) is a K-algebra which is a field but not necessarily commutative).
\end{enumerate}

